\section{Analytic models and plurivolumes}\label{sec:prelim}

\subsection{Discrepancies and the rank condition}

For a smooth model \(h:V\to Y\), log discrepancy is normalized by
\[
 a(E;Y,\Delta)
 =1+\operatorname{coeff}_E(K_{V/Y}-h^*\Delta).
\]
A \emph{sub-pair} allows negative boundary coefficients; it is
sub-log-canonical when these numbers are nonnegative.
On a smooth simple normal crossing pair, this is equivalent to every
boundary coefficient being at most one.

For comparison with the standard definition of an lc-trivial
fibration, write
\[
 A_V=K_{V/Y}-h^*\Delta,\qquad
 A_V^*=A_V+\sum_{\operatorname{coeff}_E(A_V)=-1}E.
\]
The usual rank condition is
\(\rk(f\circ h)_*\cO_V(\lceil A_V^*\rceil)=1\).
In the effective log-smooth situation of
Theorem~\ref{thm:main}, the identity is a log resolution.
For a coefficient less than one, \(\lceil-\delta\rceil=0\);
for a coefficient equal to one, the correction in \(A_Y^*\) gives
zero. Hence \(\lceil A_Y^*\rceil=0\), and connected fibers give the
rank condition. Its resolution independence is the usual
discrepancy convention. All later applications to projective
comparison families will check the generic rank directly.

\begin{lemma}\label{prelim:threshold}
Fix a model \(f':(Y',\Delta')\to X'\) as in
\eqref{eq:moduli} and a prime \(P\subset X'\).
On a joint log resolution of \(\Delta'\) and \((f')^*P\), let
\(\delta_E\) be the crepant coefficient and
\(a_E=\ord_E((f')^*P)>0\). Over a general point of \(P\),
\[
 t_P=\min_{E\mapsto P}\frac{1-\delta_E}{a_E}.
\]
In particular, the threshold is rational and unchanged by a further
resolution of the source. The divisor \(B_{X'}\) has finite support.
\end{lemma}

\begin{proof}
On the joint simple normal crossing model the vertical coefficients
of the tested sub-pair are \(\delta_E+ta_E\).
The horizontal coefficients are already at most one.
The sub-log-canonical condition is therefore precisely
\(\delta_E+ta_E\leq1\) for every vertical \(E\) dominating \(P\).
For a simple normal crossing sub-pair these bounds also test all
further exceptional valuations: the log discrepancy of a blowup
stratum is the sum of the corresponding nonnegative coordinate log
discrepancies, with any remaining smooth directions contributing
nonnegative terms. Equivalently, this is the standard local
simple normal crossing discrepancy criterion. Thus the minimum is
the infimum in \eqref{eq:threshold}.

Outside a proper analytic subset of \(X'\), the morphism and all
horizontal boundary strata are smooth and the boundary is relatively
simple normal crossing, with no vertical component. At each prime
meeting this open the same calculation gives \(t_P=1\).
Properness gives an analytic bad locus; on a compact base it has
finitely many irreducible components. Only its divisorial components
can contribute to \(B_{X'}\).
\end{proof}

\subsection{Model constructions}

We use analytic embedded and equivariant resolution, chosen projective
over the space being resolved. Functorial resolution is compatible
with smooth product directions; see
\cite[Theorems~1.1 and~1.3]{BM08}.
In particular, the resolution of a fractional monomial cover can
preserve additional smooth logarithmic coordinate directions.
Only compact or relatively compact neighborhoods are used, so no
global finiteness issue for a noncompact resolution arises.

We will also use the following standard analytic model facts.

\begin{lemma}\label{prelim:models}
\begin{enumerate}[label=\textup{(\roman*)}]
\item A closed analytic subspace of a compact K\"ahler manifold,
and a finite cover of such a space, admit smooth compact K\"ahler
models obtained by projective resolution. The corresponding
assertion holds locally over relatively compact base charts.
\item Let \(V\) be a smooth compact K\"ahler manifold and
\(V^\circ\subset V\) a dense analytic Zariski open. A finite
unramified holomorphic cover of \(V^\circ\) extends to a finite
normal cover of \(V\). After resolution it has a smooth compact
K\"ahler model unchanged over the smooth covering open.
\end{enumerate}
\end{lemma}

\begin{proof}
For \textup{(i)}, closed analytic subspaces inherit K\"ahler
structures as complex spaces. Finite maps preserve the K\"ahler
property on the source, and projective resolutions over a K\"ahler
space have K\"ahler total space. The same construction on a slightly
larger chart gives the local assertion. These are the analytic
K\"ahler model facts used in \cite[Section~4]{Fujiki78}.

For \textup{(ii)}, apply the analytic extension theorem for finite
covers, followed by normalization; a convenient formulation for
the present K\"ahler setting is
\cite[version~1, Lemma~2.21]{Villadsen24}.
Locally, after resolving the complement to normal crossings, the
cover has finite monodromy about the coordinate divisors. Sufficiently
divisible coordinate power substitutions trivialize that monodromy.
The resulting trivial finite covers extend over the polydisk, and
their finite quotients give the normal extension before the power
substitution. Normal extensions agree on overlaps by uniqueness.
Part~\textup{(i)} and projective resolution give the last assertion.
\end{proof}

A compact parameter space used below need not be K\"ahler.
Whenever a K\"ahler form on a family is needed, it will come from
one of the fixed ambient spaces of Lemma~\ref{prelim:models},
not from a presumed polarization of its parameter space.

\subsection{The prepared family}

Choose a dense analytic Zariski open \(U\subset X\) on which
\(f\), every relevant boundary stratum, and their incidence diagrams
are smooth over the base. Remove vertical boundary images.
The boundary over \(U\) is then relatively simple normal crossing.
Existence follows from properness and generic smoothness applied
to the finitely many strata.

Fix a sufficiently divisible \(m>0\) so that \(m\Delta\) is integral
and \eqref{eq:adjoint} is realized by an actual bundle isomorphism
in degree \(m\). It will be harmless to replace \(m\) by a further
multiple finitely often. Set \(d=\dim Y-\dim X\).
The line of relative log plurivolumes on \(U\) is
\begin{equation}\label{eq:open-volume-line}
 \cG_m
 =f_*\cO_{Y_U}\bigl(m(K_{Y_U/U}+\Delta|_{Y_U})\bigr)
 \simeq \cO_U\bigl(m(L-K_X)\bigr).
\end{equation}
The last notation denotes a holomorphic line bundle, whether or not
it has a global meromorphic section. The isomorphism follows from
\eqref{eq:adjoint}, the projection formula and
\(f_*\cO_Y=\cO_X\).
The evaluation map identifies its pullback with the relative
log-pluri bundle. A local generator therefore has no zeros as a
log pluriform with the specified twist; as a meromorphic relative
pluriform it has divisor \(-m\Delta\).

Take a compact K\"ahler log resolution
\(\mu_0:S\to X\), isomorphic over \(U\), such that the complement of
\(U\) is supported on a simple normal crossing divisor.
Resolve the main component of the pulled-back family to
\(p:W\to S\), unchanged over \(U\), and put
\begin{equation}\label{eq:prepared-data}
 D_W=h^*\Delta-K_{W/Y},\qquad
 L_S=\mu_0^*L,\qquad Q_S=L_S-K_S .
\end{equation}
The natural canonical-bundle identifications give
\begin{equation}\label{eq:relative-adjoint}
 K_{W/S}+D_W\sim_{\Q}p^*Q_S .
\end{equation}
We retain the actual degree-\(m\) identification coming from
\eqref{eq:adjoint}. Over \(U\), \(mQ_S\) is the line \(\cG_m\).
On \(S\), the desired moduli trace is \(M_S=Q_S-B_S\).

The local calculations below use only \eqref{eq:relative-adjoint}
and the smooth effective log-smooth generic pair.
Thus they continue to apply if the vertical coefficients of
\(D_W\) are arbitrary rational numbers.
This observation will be used for the projective comparison.

\subsection{Degenerate dimensions}

If \(X\) is a point, its rational Picard group is zero, and
Theorem~\ref{thm:main} is immediate.
If \(d=0\), connected general fibers make \(f\) bimeromorphic.
A proper bimeromorphic morphism to a smooth space is an isomorphism
outside an analytic subset of codimension at least two in the base.
At the general point of each base divisor,
\[
 t_P=1-\operatorname{coeff}_P(f_*\Delta).
\]
The actual adjoint identity on this isomorphism locus extends, with
its inverse, across codimension two. Hence
\(L=K_X+f_*\Delta\) in \(\Pic(X)_{\Q}\), and \(M_X=0\).
The same argument with the crepant sub-boundary applies on every
higher model, so all moduli traces vanish.
Here a rational divisor on a smooth base is rational Cartier;
the extension assertion is the removable-singularities theorem for
local functions expressing a line-bundle isomorphism.
We may therefore assume \(\dim X>0\) and \(d>0\) in the rest of the
proof, while allowing zero-dimensional strata and empty products.

